The results are bounded by what this study could afford, and nothing in them shows a ceiling. The classifiers are small (73 thousand parameters for the cnn and 135 thousand parameters for the rnn), each condition was trained with three seeds, and the training settings were fixed once in phase 2. With 22 test patients the 95\% intervals of the pre-registered differences are 0.06 to 0.17 wide, so a real but smaller benefit of augmentation would not be detected; not detected is not the same as absent. Changing only the input, with no retuning, raised the cnn's test rare-class macro-F1 from 0.088 to 0.371, so the models are far from saturated. Better results need more: larger models tuned on the patient-relative input, more runs per condition and a wider search of training settings, and more training patients (several databases); generators conditioned on the patient, or trained on many more patients, are the natural next test of augmentation. The answer of this study holds for its budget and should be revisited at that scale.
